% !LW recipe=lualatex x2
\documentclass[a4paper,11pt]{article}
\usepackage[margin=1.9cm,footskip=0.9cm]{geometry}
\usepackage{fontspec}
\setmainfont{TeX Gyre Heros}
\usepackage{microtype}
\usepackage{booktabs,array,tabularx,titlesec,fancyhdr}
\usepackage[table]{xcolor}
\definecolor{navy}{HTML}{172B45}
\definecolor{muted}{HTML}{536171}
\definecolor{tint}{HTML}{EFF3F7}
\setlength{\parindent}{0pt}
\setlength{\parskip}{0.58em}
\titleformat{\section}{\fontsize{12}{15}\selectfont\bfseries\color{navy}}{}{0pt}{}
\titlespacing*{\section}{0pt}{0.9em}{0.3em}
\newcolumntype{Y}{>{\raggedright\arraybackslash}X}
\pagestyle{fancy}
\fancyhf{}
\renewcommand{\headrulewidth}{0pt}
\fancyfoot[L]{\fontsize{10}{12}\selectfont\color{muted} Business Data Analytics}
\fancyfoot[R]{\fontsize{10}{12}\selectfont\color{muted}\thepage\ / 2}
\begin{document}
\raggedright
\fontsize{10.95}{14}\selectfont
{\fontsize{10}{12}\selectfont\color{muted}\textbf{EXECUTIVE BRIEF}\hfill 4 November 2026}\par
\vspace{0.18cm}
{\fontsize{25}{28}\selectfont\bfseries\color{navy}Profit-Driven Discount Strategy}\par
{\fontsize{12}{15}\selectfont\color{muted}Allocating a fixed discount budget to improve net margin}\par
\vspace{0.15cm}{\color{navy}\hrule}\vspace{0.12cm}

\section*{The decision}
Each week the retailer can send a 10\% voucher to some of its active customers. The offer can recover sales from customers who would otherwise skip a purchase. Every voucher redeemed by a customer who would have bought anyway reduces the margin on that sale.

\section*{Our recommendation}
Every Monday, rank active customers by their predicted chance of making no purchase in the next 90 days. Offer the voucher to the highest-ranked customers until adding another would exceed the expected discount budget. At about £8,200 per week, that is roughly \textbf{826 customers a week}.

\section*{Results at the same budget}
The table averages ten weekly groups from July to September 2011 that were held out from model training. Net margin is the simulated increase over 90 days for each weekly group, compared with sending no offer.

\begingroup
\setlength{\tabcolsep}{6pt}\renewcommand{\arraystretch}{1.3}
\begin{tabularx}{\textwidth}{@{}Yrrr@{}}
\toprule
\textbf{How customers are selected} & \shortstack[r]{\textbf{Expected}\\\textbf{discount spend}} & \textbf{Contacts} & \shortstack[r]{\textbf{Simulated increase}\\\textbf{in net margin}} \\
\midrule
\rowcolor{tint}\textbf{Highest risk of no purchase} & £8,213 & 826 & \textbf{£9,060} \\
Highest expected margin per contact & £8,219 & 485 & £8,228 \\
Longest time since last purchase & £8,210 & 684 & £6,079 \\
\bottomrule
\end{tabularx}\par
\endgroup

Risk ranking adds \textbf{£2,982 (49\%)} in simulated margin over the longest-inactive rule and leads in all ten test weeks under the stated assumptions.

\section*{One voucher, two outcomes}
Take a customer who usually spends £300. A 10\% voucher takes £30 off one order. At an assumed 40\% gross margin, that order leaves £120 after product costs, or £90 after the voucher. It costs £0.50 to contact the customer.

\textbf{The customer would have skipped buying.} We assume the voucher wins back 30 in every 100 of these customers. Each sale won back brings in £90 after the discount. Spread that gain across everyone contacted and deduct the cost of reaching them. The average gain per contact is \textbf{30\% × £90 − £0.50 = £26.50}.

\textbf{The customer would have bought anyway.} We assume 40 in every 100 of these customers use the voucher. The retailer makes the same sale and gives up £30 each time the voucher is used. Including the contact cost, the average effect per contact is \textbf{−40\% × £30 − £0.50 = −£12.50}.

To get the \textbf{£9,060} in the table, we repeat this calculation for each selected customer using their usual order value. Their purchase record over the following 90 days tells us which case to use. We add up the gains and losses for each weekly list and take the average across ten weeks.

\newpage
{\fontsize{21}{25}\selectfont\bfseries\color{navy}From the budget to the pilot}\par
\vspace{0.14cm}{\color{navy}\hrule}\vspace{0.1cm}

\section*{How we decide how many customers to contact}
We use the purchase history to model outcomes, but the weekly budget has to be set before those outcomes are known. For the same £30 voucher, suppose the model predicts that 60 out of 100 similar customers would skip buying. Under the assumptions on page one, 18 of those 60 would buy with a voucher. Of the other 40, 16 would use it on a purchase they already planned to make. So we expect \textbf{34 vouchers to be used}, costing £1,020 in total, or \textbf{£10.20 per offer}.

We estimate this cost for each customer, then work down the ranked list until the next offer would exceed the roughly £8,220 budget. Across ten weeks, the risk-ranked lists average \textbf{826 customers and £8,213 in expected discount spend}. We use the same budget for the rule based on time since last purchase. The 485-contact option takes up to 20\% of active customers with the highest positive expected profit per contact. The expected discount spend for that option provides the budget benchmark.

\section*{What we need to verify}
The analysis explored 5\%, 10\% and 20\% vouchers. The middle scenario is the starting point for the pilot. The records show purchases, and the 30\% recovery and 40\% redemption rates are assumptions. Product costs were also unavailable. At a 30\% gross margin, a £300 order won back would contribute £60 after the voucher. Finance should confirm actual margins before launch.

Each weekly list has a 90-day follow-up. These periods overlap, so the weekly estimates cannot be added as campaign profit.

\section*{Pilot and rollout}
With a limit of about 500 contacts per week, test the 485-contact option. In the simulation it saves 341 contacts and gives up £833 in additional margin compared with risk ranking.

\textbf{Approve a 6--8 week enrolment period.} Before launch, Finance confirms product margins and the discount cap. Marketing confirms weekly contact capacity. Analytics verifies the risk estimates on data separate from the test weeks and sets the pilot sample size.

Randomly hold back 20\% of selected customers as a no-offer control. Enrol each customer once, keep assignments fixed and follow each customer for 90 days. Compare net contribution per assigned customer after discounts, returns and contact costs. Expand if the offer produces a credible increase over the control group and discount spending stays within the approved cap.
\end{document}